\begingroup
\shorthandoff{:;!?}
\begin{tikzpicture}[
  font=\sffamily,
  fleche/.style={
    -{Stealth[length=2.6mm, width=2mm]},
    line width=0.85pt, draw=black!55
  },
  lab/.style={font=\sffamily\scriptsize, align=center, text=black!75},
  chip/.style={
    draw, line width=0.65pt, rounded corners=1.5pt,
    align=center, font=\sffamily\tiny, inner xsep=4pt, inner ysep=4pt,
    text width=2.15cm
  }
]

% --- 1. Document ---
\begin{scope}[shift={(0,0.25)}]
  \draw[line width=0.7pt, draw=black!60, fill=blue!8, rounded corners=0.6pt]
    (0.08,0.08) rectangle (1.55,1.85);
  \draw[line width=0.7pt, draw=black!70, fill=white]
    (0,0) -- (1.28,0) -- (1.47,0.22) -- (1.47,1.95) -- (0,1.95) -- cycle;
  \draw[line width=0.55pt, draw=black!70] (1.28,0) -- (1.28,0.22) -- (1.47,0.22);
  \foreach \y in {0.45,0.72,0.99,1.26,1.53}
    \draw[line width=0.45pt, draw=black!45] (0.22,\y) -- (1.22,\y);
\end{scope}
\node[lab] at (0.75,-1.05) {Document\\[-1pt]{\tiny texte + métadonnées}};

% --- 2. Segmentation (deux blocs distincts) ---
\node[chip, draw=green!55!black!55, fill=green!12] (metier) at (4.15,1.55)
  {\textbf{Segments métier}\\objet, expéditeur, type};
\node[chip, draw=violet!55, fill=violet!10] (contenu) at (4.15,0.35)
  {\textbf{Segments contenu}\\passages + chevauchement};
\node[lab] at (4.15,-1.05) {Segmentation\\[-1pt]{\tiny un document $\to$ $n$ segments}};

% --- 3. Modèle ---
\begin{scope}[shift={(7.45,0.95)}]
  \foreach \y in {0.85,0.42,0.00,-0.42}
    \fill[magenta!35] (-0.52,\y) circle (2.6pt);
  \foreach \y in {0.62,0.18,-0.28}
    \fill[magenta!50] (0.08,\y) circle (2.6pt);
  \foreach \y in {0.42,-0.05}
    \fill[magenta!65] (0.64,\y) circle (2.6pt);
  \foreach \ya in {0.85,0.42,0.00,-0.42}
    \foreach \yb in {0.62,0.18,-0.28}
      \draw[magenta!35, line width=0.35pt] (-0.52,\ya) -- (0.08,\yb);
  \foreach \ya in {0.62,0.18,-0.28}
    \foreach \yb in {0.42,-0.05}
      \draw[magenta!40, line width=0.35pt] (0.08,\ya) -- (0.64,\yb);
\end{scope}
\node[lab] at (7.50,-1.05) {Modèle d'embedding\\[-1pt]{\tiny local \textbar\ cloud}};

% --- 4. Vecteur : bande de dimensions ---
\begin{scope}[shift={(10.35,0.12)}]
  \foreach \i/\pct in {0/78, 1/22, 2/58, 3/12, 4/92, 5/38, 6/70} {
    \fill[orange!\pct, draw=orange!70!black, line width=0.35pt, rounded corners=0.6pt]
      (0, 1.68-0.28*\i) rectangle ++(0.78, 0.24);
  }
  \draw[orange!80!black, line width=1.15pt, line cap=round]
    (-0.16,1.96) -- (-0.28,1.96) -- (-0.28,-0.22) -- (-0.16,-0.22);
  \draw[orange!80!black, line width=1.15pt, line cap=round]
    (0.94,1.96) -- (1.06,1.96) -- (1.06,-0.22) -- (0.94,-0.22);
  \node[font=\sffamily\tiny, text=black!55, rotate=90] at (1.38,0.88) {dimension $d$};
\end{scope}
\node[lab] at (10.55,-1.05) {Vecteur du segment\\[-1pt]{\tiny une case $=$ une coordonnée}};

% --- 5. Espace ---
\begin{scope}[shift={(13.05,0.05)}]
  \draw[-{Stealth[length=2.2mm]}, line width=0.7pt, draw=black!70] (0,0) -- (3.05,0);
  \draw[-{Stealth[length=2.2mm]}, line width=0.7pt, draw=black!70] (0,0) -- (0,2.20);
  \foreach \x/\y in {
    0.45/0.50, 0.70/1.10, 1.05/0.38, 1.15/1.48, 1.40/0.80,
    1.55/1.82, 1.85/0.48, 1.95/1.18, 2.15/1.62, 2.35/0.70,
    2.48/1.38, 0.55/1.62, 0.90/0.75, 2.62/1.05, 1.30/1.15,
    2.15/0.28, 0.35/1.28, 2.48/1.88, 1.68/1.55, 0.80/1.88
  }
    \fill[red!55] (\x,\y) circle (1.1pt);
  \fill[green!55!black!40] (1.08,1.48) circle (1.4pt);
  \fill[green!55!black!40] (2.15,1.62) circle (1.4pt);
  \fill[violet!65] (0.70,1.10) circle (1.4pt);
  \fill[violet!65] (1.95,1.18) circle (1.4pt);
  \fill[orange!85!black] (1.58,1.02) circle (2.4pt);
  \draw[-{Stealth[length=2mm]}, line width=0.75pt, draw=black!75]
    (0,0) -- (1.58,1.02);
\end{scope}
\node[lab] at (14.55,-1.05) {Espace d'embeddings\\[-1pt]{\tiny proximité $=$ similarité}};

% Flèches : hauteurs calées entre les objets, sans traversée
\draw[fleche] (1.58,1.05) -- (3.00,1.05);
\draw[fleche] (5.32,1.05) -- (6.70,1.05);
\draw[fleche] (8.35,1.05) -- (9.85,1.05);
\draw[fleche, orange!80!black, line width=0.9pt]
  (11.55,1.05) to[out=20, in=165] (14.63,1.07);

\end{tikzpicture}
\endgroup
